% \numberwithin{equation}{section}

\chapter{Equacions MHD} \label{sec:equacionsMHD}
Les ones magnetohidrodinàmiques (MHD) són les que descriuen els esdeveniments en el plasma ionitzat de la corona solar.
Llavors s'han de tenir en compte principalment dos blocs d'equacions: les magnètiques i les de plasma.
En aquesta secció es donen les directrius principals de la deducció de les equacions MHD.
La deducció completa es troba en el Cap.~2 del llibre de \textcite{Priest_2014}, esmentat en la Introducció.

El bloc de les equacions magnètiques es conforma principalment per les equacions de~Maxwell:
\begin{gather}
	\bm{\nabla} \bm{\times} \bm{B} = \frac{1}{c^2} \frac{\partial \bm{E}}{\partial t} + \mu \bm{j}, \tag{2.1} \label{eq.2.1} \\
	\bm{\nabla} \bm{\cdot} \bm{B} = 0, \tag{2.2} \label{eq.2.2} \\
	\bm{\nabla} \bm{\times} \bm{E} = - \frac{\partial \bm{B}}{\partial t}, \tag{2.3} \label{eq.2.3} \\
	\bm{\nabla} \bm{\cdot} \bm{E} = \frac{\rho^*}{\varepsilon_0}, \tag{2.4} \label{eq.2.4}
\end{gather}
on $\bm{E}$ i $\bm{B}$ són els camps elèctric i magnètic respectivament, $\rho^*$ és la densitat de càrrega per unitat de volum,
$\bm{j}$ és la densitat de corrent, $c$ la velocitat de la llum i $\varepsilon$ i $\mu$ les permeabilitats elèctrica i magnètica del medi, respectivament.
En el cas de la corona solar el valors de les permeabilitats són propers als del buit.

Realment $\bm{B}$ és la intensitat de camp magnètic però amb el fi d'evitar una lectura feixuga es diu camp magnètic.
De totes maneres s'ha fet servir la relació constitutiva ($\bm{B} = \mu \bm{H}$), on $\bm{H}$ és directament el camp magnètic.
També es fa servir l'altra relació constitutiva ($\bm{D} = \varepsilon \bm{E}$), on $\bm{E}$ és el camp elèctric i $\bm{D}$ el vector desplaçament.

Per seguir amb les equacions magnètiques es té en compte que la pertorbació d'una protuberància quiescient suposa un problema no relativista.
Llavors la llei d'Ohm dicta que el corrent és proporcional al camp elèctric total:
\begin{equation}
	\bm{j} =  \sigma (\bm{E} + \bm{v} \bm{\times} \bm{B}), \tag{2.5} \label{eq.2.5}
\end{equation}
i si s'aïlla el camp elèctric estàtic:
\begin{equation}
	\bm{E} = - \bm{v} \bm{\times} \bm{B} + \bm{j}/\sigma. \tag{2.6} \label{eq.2.6}
\end{equation}

Finalment s'obté l'equació d'inducció al substituir el camp elèctric estàtic %de l'Eq.~(\ref{eq.2.6})
dins l'Eq.~(\ref{eq.2.3}):
\begin{equation}
	\frac{\partial \bm{B}}{\partial t} = \bm{\nabla} \bm{\times} (\bm{v} \bm{\times} \bm{B}). \tag{2.7} \label{eq.2.7}
\end{equation}
És important tenir en compte que es fa la suposició de MHD ideal, llavors es considera una resistivitat del medi, $\sigma$,
infinita, i el terme $\bm{j}/\sigma$ desapareix de l'Eq.~(\ref{eq.2.6}).
En aquest cas és possible fer l'assumpció ja que les escales de longitud del problema són grans.
Tot i així, en algunes regions d'alta concentració de corrent o a escales espacials petites no sempre es pot fer aquesta aproximació.

Les variables més importants en un plasma són el camp magnètic, $\bm{B}$, i la velocitat del propi plasma, $\bm{v}$, per tant, és convenient escriure totes les expressions en funció d'aquestes.
En canvi, les variables de corrent, $\bm{j}$, i de camp elèctric, $\bm{E}$, són secundàries, de manera que poden ser substituides per les principals.
Amb aquesta consideració es simplifiquen les equacions ja que apareixen manco variables, fet que s'agraeix a l'hora d'implementar i resoldre les equacions amb l'ajuda de Dedalus, com es veurà més envant.



El segon bloc d'equacions a tractar és el de les equacions del plasma. Les tres equacions principals són les que segueixen, començant per la de continuïtat de massa:
\begin{equation}
	\frac{d\rho}{dt} + \rho \bm{\nabla} \cdot \bm{v} = 0 \quad (\text{a}) \quad , \quad \frac{\partial \rho}{\partial t} + \nabla \cdot (\rho \bm{v}) = 0 \quad (\text{b}), \tag{2.8} \label{eq.2.8}
\end{equation}
on s'ha definit la derivada material o convectiva com
\begin{equation}
	\frac{d}{dt} = \frac{\partial}{\partial t} + \bm{v} \cdot \bm{\nabla}. \tag{2.9} \label{eq.2.9}
\end{equation}

La segona equació rellevant realment és el subconjunt de les equacions del moviment del plasma sota neutralitat elèctrica:
\begin{equation}
	\rho \frac{d \bm{v}}{dt} = - \bm{\nabla} p + \bm{j} \bm{\times} \bm{B} + \bm{F}, \tag{2.10} \label{eq.2.10}
\end{equation}
on no es tenen en compte els efectes incercials i les forces ($\bm{F}$) poden ser viscoses ($\bm{F}_{v}$) o gravitatòries ($\bm{F}_g$).
D'aquesta manera, al resoldre aquestes equacions s'obtenen ones magneto-viscoses o magneto-gravitacionals respectivament.
No obstant això, en el cas d'aquest estudi es suposen forces externes nul·les ($\bm{F} = \bm{0}$) per les qüestions que s'han aclarit en la introducció (\hyperref[sec:intro]{Sec. 1}). %(\hyperref[sec:intro]{\nameref{sec:intro}}). %
Per consegüent les ones resultants són magneto-acústiques (també anomenades magnetohidrodinàmiques, MHD), que corresponen a l'Eq.~(\ref{eq.2.10}) si només es té en compte la força de Lorentz ($\bm{j} \bm{\times} \bm{B}$) i el gradient de pressió ($- \bm{\nabla} p$).

Finalment la tercera equació de plasma rellevant és la de gas ideal, que pot ser escrita de diverses maneres:
\begin{equation}
	p = \frac{\tilde{R}}{\tilde{\mu}} \rho T = \frac{k_B}{m} \rho T = n k_B T, \tag{2.11} \label{eq.2.11}
\end{equation}
on $k_B$ és la constant de Boltzmann, $m_p$ és la massa del protó, $\tilde{R} = k_B/m_p$ la constant dels gasos ideals, $\tilde{\mu} = m/m_p$ el pes atòmic
mitjà, $\rho = \tilde{\mu} m_p n = m n$ la densitat, $m$ la massa de partícula mitja i $n$ el nombre total de partícules.


Una hipòtesi important amb profundes implicacions que es fa en aquest estudi és la d'adiabaticitat. La conservació d'energia del sistema modifica l'equació de calor
\begin{equation}
	\rho T \frac{ds}{dr} = - \mathcal{L}, \tag{2.12} \label{eq.2.12}
\end{equation}
ò alternativament
\begin{equation}
	\frac{\rho^\gamma}{\gamma - 1} \frac{d}{dt} \left( \frac{p}{\rho^\gamma} \right) = - \mathcal{L}, \tag{2.13} \label{eq.2.13}
\end{equation}
de manera que
\begin{equation}
	\frac{d}{dt} \left( \frac{p}{\rho^\gamma} \right) = 0, \tag{2.14} \label{eq.2.14}
\end{equation}
on $s$ és l'entropia, $\gamma = c_p/c_v = 5/3$ la proporció de calors específics i $\mathcal{L}=0$ és la funció pèrdua d'energia, nul·la en el cas adiabàtic.
Així doncs, de l'equació de calor s'obté una condició que més envant es farà servir al resoldre les equacions.


Per acabar aquesta secció es reescriuen totes les assumpcions fetes en la deducció de les equacions.
Aquestes són extretes del Cap.~2 del llibre de \textcite{Priest_2014}.
\begin{enumerate}[label=\Roman*.]
	\item \label{punt:I} Plasma continu i en equilibri estàtic.
	\item \label{punt:II} Plasma en equilibri termodinàmic (TD).
	\item \label{punt:III} Coeficient $\sigma \rightarrow \infty \implies \eta \rightarrow 0$ i $\mu$ constant, propietats del plasma isotròpiques.
	\item \label{punt:IV} Equacions en sistema de referència del plasma, no del Sol.
	\item \label{punt:V} Efectes relativistes ignorats, velocitats molt menors a la del llum.
	\item \label{punt:VI} Plasma fluid únic, tot i que existeixen models de 2 i 3.
	\item \label{punt:VII} Llei d'Ohm simplificada.
\end{enumerate}